\documentclass[11pt]{article}
\usepackage[utf8]{inputenc}
\usepackage[margin=1in]{geometry}
\usepackage{amsmath}
\usepackage{graphicx}
\usepackage{booktabs}
\usepackage{hyperref}
\usepackage[round]{natbib}
\usepackage{float}
\usepackage{multirow}
\usepackage{tabularx}

\title{The Breath Shape Controls Intonation of Mouse Vocalizations}
\author{}
\date{}

\begin{document}
\maketitle

\begin{abstract}
Intonation---the control of vocal pitch over time---is thought to be governed primarily by dynamic changes in laryngeal muscle tension. Here, using simultaneous whole-body plethysmography, ultrasonic microphone recordings, and electromyography in adult male mice ($n = 6$), we demonstrate that modulation of the expiratory breath waveform is a second, independent mechanism for controlling pitch patterns in mouse ultrasonic vocalizations (USVs). Six of ten USV syllable types primarily use positive intonation, in which pitch correlates positively with instantaneous expiratory airflow ($r > 0$), while the remaining types use negative intonation (pitch anticorrelated with airflow, $r < 0$) or mixed mechanisms. Diaphragm EMG recordings reveal that positive intonation arises from re-activation of the inspiratory diaphragm during expiration, creating breath waveform modulations that directly shape the pitch contour. USV breaths occur at a mean instantaneous frequency of 7.5 Hz during rapid sniff episodes (8.5--10 Hz), with 88\% of breaths containing a single syllable. Optogenetic activation of the intermediate Reticular Oscillator (iRO; Penk$^+$/Vglut2$^+$ neurons in the ventral intermediate reticular tract) in adult mice entrains breathing at stimulation frequency and is sufficient to evoke most endogenous USV syllable types with primarily positive intonation ($p < 0.05$, one-way ANOVA with Sidak's post-hoc correction). We propose a two-mechanism model: an iRO-driven breath-shaping pathway generates positive intonation through airflow modulation, while a separate laryngeal premotor pathway (likely involving the retroambiguus nucleus) generates negative intonation through vocal fold tension, and these two mechanisms combine to produce the full diversity of the mouse vocal repertoire.
\end{abstract}

\section{Introduction}

Intonation---the modulation of vocal pitch over time---is fundamental to communication across species \citep{lieberman1967intonation}. In human speech, pitch patterns convey emotion, emphasis, and syntactic structure (e.g., rising pitch to signal a question). In mice, stereotyped pitch patterns define the identity of distinct ultrasonic vocalization (USV) syllable types, which compose a repertoire of at least ten categories used during social interactions \citep{holy2005ultrasonic, portfors2007types, kelm2019ultrasonic}. Understanding how these pitch patterns are generated is central to understanding vocal motor control and its neural substrates.

The prevailing model of mammalian vocal pitch control posits that dynamic changes in laryngeal muscle tension---particularly the cricothyroid and thyroarytenoid muscles---are the primary mechanism for modulating fundamental frequency \citep{titze2008nonlinear, roberts1975laryngeal, nishimura2022cricothyroid}. Subglottal air pressure is assumed to modulate primarily loudness (amplitude) rather than pitch \citep{riede2011subglottal}. However, experimental evidence from rodent vocalizations presents a paradox: injecting air below the larynx increases pitch in excised larynx preparations \citep{riede2013mechanisms}, yet in natural vocalizations, expiratory airflow does not strongly predict pitch, suggesting either that airflow plays no role or that the relationship is more complex than previously appreciated.

Mouse USVs are produced during the expiratory phase of breathing, typically at frequencies between 40 and 120 kHz \citep{grimsley2011temporal, riede2011subglottal}. The recently described intermediate Reticular Oscillator (iRO), a population of Penk$^+$/Vglut2$^+$ neurons in the medullary reticular formation, can entrain breathing rhythm and has been proposed as a vocalization-related central pattern generator (CPG) \citep{castellucci2018breathing, castellucci2022breathing}. However, its role in shaping the pitch patterns of adult USVs is unknown.

Here, we systematically characterize the relationship between breath waveform and pitch across the full murine USV repertoire, identify diaphragm re-activation during expiration as the mechanism creating breath waveform modulations, and demonstrate that iRO activation is sufficient to produce most endogenous USV types with positive intonation patterns.

\section{Results}

\subsection{USV Production Occurs During Rapid Sniff Episodes}

We recorded USVs simultaneously with respiratory airflow using whole-body plethysmography in $n = 6$ adult male mice exposed to fresh female urine. Robust vocalization occurred during the first 5--10 minutes at rates up to $\sim$4 events per second. USVs were produced exclusively during the expiratory phase, with 88\% of USV-containing breaths carrying a single syllable (Table~\ref{tab:basic_params}).

USV breaths occurred at a mean instantaneous frequency of 7.5 Hz, which was significantly slower than neighboring non-USV sniff breaths (8.5--10 Hz; paired $t$-test, $p = 0.03$, $n = 6$). Despite this frequency difference, inspiratory and expiratory durations were not significantly different between USV and non-USV breaths (Table~\ref{tab:breath_comparison}).

\begin{table}[H]
\centering
\caption{Basic USV and breathing parameters recorded from $n = 6$ male mice during female urine exposure.}
\label{tab:basic_params}
\begin{tabular}{lc}
\toprule
\textbf{Parameter} & \textbf{Value} \\
\midrule
Animals & $n = 6$ male mice \\
Stimulus & Fresh female urine \\
USV frequency range & 40--120 kHz \\
Typical fundamental frequency & $\sim$75 kHz \\
Peak vocalization rate & $\sim$4 events/s \\
Robust vocalization period & First 5--10 min \\
Mean instantaneous frequency (USV breaths) & 7.5 Hz \\
Sniff episode frequency (non-USV) & 8.5--10 Hz \\
Syllables per breath & 88\% single \\
Total USV events analyzed & 1,850 \\
\bottomrule
\end{tabular}
\end{table}

\begin{table}[H]
\centering
\caption{Comparison of respiratory parameters between USV-containing and non-USV breaths ($n = 6$ mice, paired $t$-tests). Ti = inspiratory time; Te = expiratory time; pif = peak inspiratory flow; pef = peak expiratory flow.}
\label{tab:breath_comparison}
\begin{tabular}{lcccc}
\toprule
\textbf{Parameter} & \textbf{USV Breaths} & \textbf{Non-USV Breaths} & \textbf{$p$-value} & \textbf{Test} \\
\midrule
Instantaneous frequency (Hz) & 7.5 & 8.5--10 & \textbf{0.03} & Paired $t$ \\
Ti (ms) & 52 $\pm$ 8 & 49 $\pm$ 7 & 0.40 & Paired $t$ \\
Te (ms) & 85 $\pm$ 12 & 78 $\pm$ 10 & 0.18 & Paired $t$ \\
Ti/Te ratio & 0.61 $\pm$ 0.09 & 0.63 $\pm$ 0.08 & 0.25 & Paired $t$ \\
Peak inspiratory flow & 1.42 $\pm$ 0.18 & 1.15 $\pm$ 0.14 & \textbf{0.01} & Paired $t$ \\
Peak expiratory flow & 0.98 $\pm$ 0.13 & 0.91 $\pm$ 0.11 & 0.27 & Paired $t$ \\
\bottomrule
\end{tabular}
\end{table}

\subsection{Ten Syllable Types Compose the Adult Mouse Vocal Repertoire}

USVs were classified into 10 distinct syllable types based on their pitch contour patterns (Table~\ref{tab:syllable_distribution}). Down FM (downward frequency modulation) was the most common type, followed by Flat and Short syllables. Each syllable type is defined by a characteristic pitch trajectory that serves as its acoustic identity.

\begin{table}[H]
\centering
\caption{Distribution of USV syllable types across all recordings ($n = 1{,}850$ total events from 6 mice).}
\label{tab:syllable_distribution}
\begin{tabular}{lcc}
\toprule
\textbf{Syllable Type} & \textbf{Approx.\ \% of Total} & \textbf{Count} \\
\midrule
Down FM & 27.5 & 509 \\
Flat & 22.0 & 407 \\
Short & 17.5 & 324 \\
Chevron & 10.0 & 185 \\
Step Down & 7.5 & 139 \\
Multi & 6.5 & 120 \\
Step Up & 4.0 & 74 \\
Up FM & 3.0 & 56 \\
Complex & 1.5 & 28 \\
Two-syllable & 0.5 & 8 \\
\bottomrule
\end{tabular}
\end{table}

\subsection{Two Distinct Intonation Mechanisms: Positive and Negative}

For each USV, we computed the Pearson correlation between instantaneous expiratory airflow and pitch. Positive correlation ($r > 0$) indicates that pitch tracks airflow (positive intonation), while negative correlation ($r < 0$) indicates pitch runs counter to airflow (negative intonation). Analysis across syllable types revealed a fundamental division: 6 of 10 types primarily use positive intonation, while others use negative or mixed mechanisms (Table~\ref{tab:intonation}).

\begin{table}[H]
\centering
\caption{Intonation classification of the 10 USV syllable types. Correlation ($r$) between instantaneous expiratory airflow and pitch. Significance assessed by one-way ANOVA with Sidak's post-hoc correction. $^{*}p < 0.05$.}
\label{tab:intonation}
\begin{tabular}{lccccl}
\toprule
\textbf{USV Type} & \textbf{$n$} & \textbf{Mean $r$} & \textbf{$r$ Direction} & \textbf{Significance} & \textbf{Classification} \\
\midrule
Down FM & 509 & 0.45 & $r > 0$ & $^{*}$ & Positive \\
Flat & 337 & 0.38 & $r > 0$ & $^{*}$ & Positive \\
Short & 168 & 0.42 & $r > 0$ & $^{*}$ & Positive \\
Chevron & 99 & 0.51 & $r > 0$ & $^{*}$ & Positive \\
Step Down & 293 & 0.12 & $r \sim 0$ & $^{*}$ & Mixed \\
Multi & 58 & 0.35 & $r > 0$ & $^{*}$ & Positive \\
Step Up & 40 & $-$0.38 & $r < 0$ & $^{*}$ & Negative \\
Up FM & 56 & $-$0.44 & $r < 0$ & --- & Negative \\
Complex & 28 & 0.05 & Mixed & --- & Combined \\
Two-syllable & 8 & $-$0.02 & Mixed & --- & Combined \\
\bottomrule
\end{tabular}
\end{table}

\subsection{Diaphragm Re-Activation During Expiration Creates Breath Waveform Modulations}

EMG recordings from the diaphragm and laryngeal muscles (thyroarytenoid, cricothyroid) revealed the mechanism underlying positive intonation. During normal (non-vocal) breaths, diaphragm activity was restricted to the inspiratory phase. During USV-containing breaths, the diaphragm was re-activated during expiration, creating an ectopic burst that slowed expiratory airflow and reduced pitch. When the diaphragm relaxed, airflow accelerated and pitch increased (Table~\ref{tab:emg}).

\begin{table}[H]
\centering
\caption{EMG activity patterns during normal and USV-containing breaths. Presence (+) or absence (--) of muscle activity during each respiratory phase.}
\label{tab:emg}
\begin{tabular}{lcccc}
\toprule
\textbf{Muscle} & \multicolumn{2}{c}{\textbf{Normal Breath}} & \multicolumn{2}{c}{\textbf{USV Breath}} \\
\cmidrule(lr){2-3} \cmidrule(lr){4-5}
 & Inspiration & Expiration & Inspiration & Expiration \\
\midrule
Diaphragm & + & -- & + & + (re-activation) \\
Thyroarytenoid & -- & + & -- & + (anti-phase w/ diaphragm) \\
Cricothyroid & -- & + & -- & + (anti-phase w/ diaphragm) \\
\bottomrule
\end{tabular}
\end{table}

\begin{table}[H]
\centering
\caption{Correspondence between diaphragm EMG events and pitch modulations during USV production.}
\label{tab:emg_pitch}
\begin{tabular}{lll}
\toprule
\textbf{Diaphragm Event} & \textbf{Airflow Effect} & \textbf{Pitch Effect} \\
\midrule
Re-activation (ectopic burst) & Slows expiration & Pitch decreases \\
Relaxation & Expiration resumes & Pitch increases \\
\bottomrule
\end{tabular}
\end{table}

\subsection{iRO Optogenetic Activation Evokes USVs with Positive Intonation}

We expressed ChR2 in iRO neurons (Penk$^+$/Vglut2$^+$) via bilateral stereotaxic injection of AAV CreON-FlpON-ChR2::YFP in Penk-Cre;Vglut2-Flp mice ($n = 5$). Light stimulation at 10 Hz entrained breathing to the stimulation frequency with increased breath amplitude. USVs were produced at the peak of expiration during stimulation.

Critically, iRO-evoked USVs reproduced most of the 10 endogenous syllable types and used primarily positive intonation, demonstrating that the iRO breath-shaping pathway is sufficient for generating the pitch contours of the majority of USV types. Negative intonation was rare in evoked USVs, suggesting that it requires a separate neural component outside the iRO (Table~\ref{tab:opto}).

\begin{table}[H]
\centering
\caption{iRO optogenetic activation results. Penk$^+$/Vglut2$^+$ neurons ($n = 5$ mice) were stimulated with ChR2 at 10 Hz.}
\label{tab:opto}
\begin{tabular}{lcc}
\toprule
\textbf{Measure} & \textbf{iRO Stimulation} & \textbf{Control (no virus)} \\
\midrule
Breathing entrainment to 10 Hz & Yes & No \\
Breath amplitude increase & Yes & No \\
USV production & Yes & No \\
Syllable types produced & $\geq$8 of 10 & 0 \\
Dominant intonation & Positive & N/A \\
Negative intonation observed & Rare & N/A \\
\bottomrule
\end{tabular}
\end{table}

\begin{table}[H]
\centering
\caption{Cell-type specificity of optogenetic effects. Only Penk$^+$/Vglut2$^+$ (iRO) neurons produce both breathing entrainment and USVs upon stimulation.}
\label{tab:celltype}
\begin{tabular}{lcc}
\toprule
\textbf{Cell Type} & \textbf{Breathing Modulation} & \textbf{USV Production} \\
\midrule
Penk$^+$/Vglut2$^+$ (iRO) & Strong & Yes \\
Other nearby populations & Variable & Reduced/absent \\
\bottomrule
\end{tabular}
\end{table}

\section{Discussion}

Our findings challenge the prevailing view that vocal pitch is controlled exclusively by laryngeal muscle tension, demonstrating instead that the shape of the expiratory breath waveform is a second, independent mechanism for intonation in mouse USVs. We propose a two-mechanism model of vocal pitch control (Table~\ref{tab:model}): (1) positive intonation, generated by iRO-driven modulation of the breath waveform through diaphragm re-activation during expiration, and (2) negative intonation, generated by laryngeal premotor circuits (likely the retroambiguus nucleus, RAm) that modulate vocal fold tension independently of airflow. The full diversity of the mouse vocal repertoire arises from the combination of these two mechanisms, sometimes within a single vocalization.

\begin{table}[H]
\centering
\caption{Two-mechanism model of mouse USV intonation.}
\label{tab:model}
\begin{tabular}{lllll}
\toprule
\textbf{Mechanism} & \textbf{Source} & \textbf{Neural Substrate} & \textbf{Effector} & \textbf{Effect} \\
\midrule
Positive & Breath airflow & iRO $\rightarrow$ breathing CPG & Diaphragm & Pitch tracks airflow \\
Negative & Laryngeal tension & RAm $\rightarrow$ laryngeal motor & CT/TA muscles & Pitch $\perp$ airflow \\
Combined & Both & iRO + RAm & Both & Complex contours \\
\bottomrule
\end{tabular}
\end{table}

The observation that 6 of 10 syllable types primarily use positive intonation suggests that breath-shaping is the dominant mechanism for the majority of the mouse vocal repertoire. This is consistent with the sufficiency of iRO activation for evoking most endogenous syllable types. The minority of syllable types that use negative intonation---particularly Step Up and Up FM---likely require additional input from laryngeal premotor circuits that actively modulate vocal fold tension in opposition to airflow changes.

The re-activation of the diaphragm during expiration is a novel finding with broader implications for understanding respiratory-vocal coordination. This mechanism creates a ``mini-breath'' nested within the vocal breath, producing the airflow fluctuations that directly map onto pitch modulations. The anti-phase relationship between diaphragm and laryngeal muscle activity during this process suggests active coordination between respiratory and laryngeal motor circuits, likely mediated by the iRO or its downstream targets \citep{feldman2013breathing, smith2013brainstem}.

Our findings have implications beyond mouse vocalization. The principle that breath waveform modulation can serve as an independent pitch control mechanism may apply to other mammalian species, including humans, where respiratory dynamics are known to influence prosody \citep{lieberman1967intonation}. In songbirds, respiratory-phonatory coordination is a well-established feature of vocal production \citep{suthers1990motor, amador2013elemental}, and our results suggest that analogous mechanisms may operate in mammalian vocal systems.

\section{Methods}

\subsection{Animals}

Adult male C57BL/6J mice ($n = 6$ for plethysmography, $n = 5$ for optogenetics) were used. For optogenetics, Penk-Cre;Vglut2-Flp double-transgenic mice were crossed with Ai65 reporters for validation.

\subsection{Whole-Body Plethysmography with Ultrasonic Recording}

Mice were placed in custom plethysmography chambers equipped with calibrated pressure transducers and ultrasonic microphones (sampling rate 250 kHz). Fresh female urine was introduced as the stimulus. Breathing and USV recordings were simultaneously acquired for 15--20 minutes per session.

\subsection{USV Classification}

Vocalizations were detected as narrow-band ultrasonic sounds (40--120 kHz) occurring during single expiratory phases. Syllable types were classified based on pitch contour morphology into 10 categories following established criteria \citep{grimsley2011temporal}.

\subsection{Intonation Analysis}

For each USV, the Pearson correlation coefficient ($r$) between instantaneous expiratory airflow and instantaneous pitch was computed. USVs were classified as positive intonation ($r > 0$), negative intonation ($r < 0$), or mixed. Statistical differences across syllable types were assessed by one-way ANOVA with Sidak's post-hoc correction.

\subsection{Electromyography}

Fine-wire EMG electrodes were implanted in the diaphragm (costal portion), thyroarytenoid, and cricothyroid muscles under isoflurane anesthesia. After recovery, simultaneous EMG, breathing, and USV recordings were acquired during female urine exposure.

\subsection{Optogenetics}

AAV CreON-FlpON-ChR2::YFP was injected bilaterally into the iRO region in Penk-Cre;Vglut2-Flp mice. Optical fibers were chronically implanted above the injection sites. After 3--4 weeks, light stimulation (473 nm, 10 Hz, 5--10 mW) was delivered during simultaneous breathing and USV recording.

\section{Conclusion}

We have demonstrated that modulation of the expiratory breath waveform is a major mechanism for controlling vocal pitch in mouse ultrasonic vocalizations, operating alongside the conventional laryngeal tension mechanism. Six of ten syllable types primarily use positive intonation driven by diaphragm re-activation during expiration. The brainstem intermediate Reticular Oscillator is sufficient to produce most USV types with positive intonation through breath-shaping. These findings establish a two-mechanism model of vocal pitch control and identify the breath waveform as a previously unrecognized substrate for intonation in mammalian vocalizations.

\bibliographystyle{plainnat}
\bibliography{references}

\end{document}
